\section{Inversions of Test Data}%
\label{sec:InversionsofTestData}%

\begin{outline}
\item The test data are simulated \ESISII\ observations of scenes whose line profiles are known,
    so the inversions can be compared with the truth.
\end{outline}%
\subsection{Test Data}%
\label{subsec:TestData}%

\begin{outline}
\item Scenes with known line profiles: \IRIS\ \SiIV\ rasters of \ARs\ rescaled to
    \NeVII~\citep{DePontieu2014}; \AIA\ images given the radiance and width of
    \SiXII; and beam heating events from the HYDRAD simulations inserted into the moss.
    The first two are \texttt{esis.data.synth.scene\_iris} and \texttt{scene\_aia}.
\item A radiative \MHD\ simulation of an \AR\ with MURaM~\citep{Rempel2023}, whose
    velocity-resolved \DEM\ gives the line profiles directly, with
    \texttt{esis.flights.f2.data.synth.scene\_muram} (Figure~\ref{fig:muram}).
\item Forward modeled through \texttt{esis.flights.f2}, with photon and read noise for a single
    exposure.
\item \textbf{Figure:} a simulated image from every channel.
\end{outline}%


\begin{figure*}[!htb]%
\centering%
\includegraphics{generated/muram.pdf}%
\caption{
Above, a radiative \MHD\ simulation of an active region with MURaM, before its
M6.6 flare~\citep{Rempel2023}: the radiance and the Doppler centroid of \NeVII\ and
\SiXII, positive away from the observer, synthesized from its velocity-resolved
\DEM.
The plasma in each bin of temperature and of line-of-sight velocity emits a Gaussian
at the center of the bin, with its thermal width and the width of plasma spread
evenly across the bin.
Each pixel of the simulation is stretched to span one pixel of the detector, so that
the simulation fills the field of view at \stretchMuram\ times its true size.
The Doppler centroid is not drawn, in gray, where the radiance is among the faintest
\fractionMaskedMuram\ of its map.
Below, a simulated \timeExposure\ exposure from one channel of \ESISII\ of the
\numLinesMuram\ lines of the passband which carry at least \fractionLineMin\ of its
radiance, in the electrons collected by each pixel of the \CCD, on a square-root
scale, with the shot noise of the photons and the read noise of the camera.}%
\label{fig:muram}%
\end{figure*}

%
\subsection{Inversion Method}%
\label{subsec:InversionMethod}%

\begin{outline}
\item \MARTCapital~\citep{Gordon1970}, as applied to \ESISI~\citep[][Appendix~B]{Parker2022}:
    multiplicative updates keep the solution positive and assume no line shape, and regularization
    suppresses the ``plaid'' artifacts along the projection directions.
\item Implemented in \texttt{ctis}.
\end{outline}

%
\subsection{Validation Experiments}%
\label{subsec:ValidationExperiments}%

\begin{outline}
\item Invert using every subset of three or more channels, and compare the recovered moments with
    each other and with the truth.
\item Invert using a subset, forward model the withheld channels, and compare them with the
    simulated observations.
    Only this experiment can be repeated with the flight data, so the test data calibrate how well
    the withheld-channel residual predicts the true error.
\item Error in intensity, Doppler shift, and width versus the number of channels, separately for
    compact brightenings in \NeVII\ and extended loops in \SiXII.
\end{outline}

%
\subsection{Results}%
\label{subsec:Results}%

\begin{outline}
\item Compact \NeVII\ brightenings: whether the blue shifts of beam heating events are recovered
    with few channels, as for the explosive events of \ESISI~\citep{Parker2022}.
\item Extended \SiXII\ structures: \numChannelsRequired\ channels recover the Doppler shift to
    within \velocityError.
\item \textbf{Figures:} recovered versus true moments; error versus the number of channels.
\end{outline}
